% =====================================================================
%  Front matter.
%  Source: tier0.md:1-121 (title blurb, status, how-to-study, contents).
%  The <details> "full section list" (tier0.md:64-119) is not dropped --
%  it is reproduced verbatim as Appendix A.
% =====================================================================

\begin{titlepage}
\centering
\vspace*{3.5cm}

{\Huge\bfseries Tier 0 --- Substrate\par}

\vspace{1.2em}
{\Large Physical foundations, the reaction network as a formal object,\\[0.3em]
the dataset, and the modelling frame\par}

\vspace{0.8em}
{\large\itshape A study report for the conservation-by-construction GNN\\
emulator of silicon burning\par}

\vspace{3em}
\begin{minipage}{0.82\textwidth}
\small
Complete deep dive: physical foundations, project motivation, S0 (species and
abundance algebra), S1 (the dataset), the modelling frame, and the numerical
and ML prerequisites. Everything needed before Tier 1 (reaction rates and
detailed balance), in one document.
\end{minipage}

\vfill
{\small Literature audit 2026-08-12\par}
\vspace{0.5em}
{\small Source document: \code{study/tier0.md} \quad$\cdot$\quad
Audit log: \code{study/tier0-audit.md}\par}
\end{titlepage}

% ---------------------------------------------------------------------
\chapter*{Status of this document}
\addcontentsline{toc}{chapter}{Status of this document}
\label{sec:status}

\begin{warn}
\textbf{Personal study material, not project spec.} Deliberately written
without reference to \code{docs/} --- everything is derived from first
principles or read off the business code, configs, notebooks, and tests.
\end{warn}

\begin{itemize}
\item Textbook physics and derivations here are \textbf{mine to check}, not
      citable project output.

\item Measured project numbers quoted here originate in \code{RESULTS.md} with
      script $+$ commit $+$ data provenance. This document is \textbf{not}
      their source of truth; if a number here disagrees with \code{RESULTS.md},
      \code{RESULTS.md} wins.

\item \textbf{Literature audit 2026-08-12.} Every scholarly claim, constant,
      and derivation chain was verified against primary sources (arXiv \TeX{}
      cached in \code{data/literature/}, AME2020/ENSDF/CODATA data pages) with
      all worked arithmetic recomputed in float64; errors found were corrected
      in place. Inline author--year citations point to the References section
      at the end; the claim-by-claim audit log, including every correction with
      its before/after and supporting source passage, is
      \code{study/tier0-audit.md}.
\end{itemize}

Notation follows the repo: unicode math, \nuc{Co}{55} in prose, \code{co55} in
code-adjacent contexts, \Ye{} / $\nu$ / $\varphi$ / $\kappa$ as in the
codebase.

\paragraph{A note on this edition.} This report re-architects the source
document's Part~0--VII ordering into a narrative sequence: astrophysical
context, then the physical substrate, then the network as a formal object,
then the computational problem, the dataset, the modelling frame, numerical
analysis, ML framing, and working practice. \emph{No content has been
omitted.} Appendix~\ref{app:sectionmap} reproduces the source document's own
section map verbatim, and Appendix~\ref{app:concordance} is a
heading-by-heading concordance from the original numbering to this one,
including the three self-check items that were relocated so that no question
tests material the reader has not yet reached.

% ---------------------------------------------------------------------
\chapter*{How to study each node}
\addcontentsline{toc}{chapter}{How to study each node}
\label{sec:howtostudy}

\begin{figure}[H]
\centering
\begin{tikzpicture}[
  font=\small,
  stepbox/.style={draw, rounded corners=2pt, line width=0.7pt,
               minimum height=8mm, inner xsep=6pt, fill=resultrule!6},
  num/.style={circle, draw, line width=0.7pt, inner sep=1.2pt,
              minimum size=5.2mm, font=\footnotesize\bfseries,
              fill=white}]

\node[stepbox] (a) {\textbf{DERIVE}};
\node[stepbox, below=5mm of a] (b) {\textbf{READ}};
\node[stepbox, below=5mm of b] (c) {\textbf{ORACLE}};
\node[stepbox, below=5mm of c] (d) {\textbf{SEE}};

% align the boxes to a common left edge and equal width
\begin{scope}[on background layer]
\node[fit=(a)(b)(c)(d), inner sep=0pt] (grp) {};
\end{scope}

\node[num] at ([xshift=-9mm]a.west) {1};
\node[num] at ([xshift=-9mm]b.west) {2};
\node[num] at ([xshift=-9mm]c.west) {3};
\node[num] at ([xshift=-9mm]d.west) {4};

\node[anchor=west, text width=0.66\textwidth] at ([xshift=6mm]a.east)
  {pen-and-paper: get the equation from physics before reading code};
\node[anchor=west, text width=0.66\textwidth] at ([xshift=6mm]b.east)
  {the module, top-to-bottom; its docstring states the contract};
\node[anchor=west, text width=0.66\textwidth] at ([xshift=6mm]c.east)
  {the test file --- it encodes what ``correct'' means, numerically};
\node[anchor=west, text width=0.66\textwidth] at ([xshift=6mm]d.east)
  {the notebook figure --- the measured behaviour on real data};

\draw[-{Latex[length=2mm]}, line width=0.6pt] (a) -- (b);
\draw[-{Latex[length=2mm]}, line width=0.6pt] (b) -- (c);
\draw[-{Latex[length=2mm]}, line width=0.6pt] (c) -- (d);
\end{tikzpicture}
\caption{The four-step study loop for each node of the curriculum.}
\label{fig:studyloop}
\end{figure}

Tests are the physics oracle here (project rule: \emph{tests before
implementation for anything with a physics oracle}), and module docstrings
state \emph{contracts} rather than describing code. So step~3 is where
correctness is pinned.

% ---------------------------------------------------------------------
\chapter*{Reading map}
\addcontentsline{toc}{chapter}{Reading map}
\label{sec:readingmap}

\begin{tabularx}{\textwidth}{@{}L{2.4cm} L{3.6cm} Y@{}}
\toprule
\textbf{Part} & \textbf{Theme} & \textbf{Covers} \\
\midrule
\ref{part:context} & Astrophysical context &
  Stellar context, the burning sequence, \Ye{} and explodability, the
  observational grounding \\
\ref{part:substrate} & The physical substrate &
  Plasma thermodynamics \& EOS, statistical mechanics, nuclear structure \&
  data, neutrinos \\
\ref{part:network} & The network as a formal object &
  $Y = X/A$, \Ye{}, the two networks, the bipartite graph, reaction
  nomenclature, Si-burning topology, the timescale hierarchy \\
\ref{part:computational} & The computational problem &
  The cost problem, softwired networks, why ML emulation, conservation, why a
  GNN \\
\ref{part:dataset} & The dataset &
  Regime box, $\dd t$ grid, Sobol design, identity, label noise, splits \\
\ref{part:frame} & The modelling frame &
  Operator splitting, one-zone \& the reachable manifold, semigroup \& rollout \\
\ref{part:numerics} & Numerical analysis &
  Catastrophic cancellation ($=\kappa$), conditioning, stiffness, floating
  point, Newton \& globalization, ODE error control \\
\ref{part:ml} & ML framing &
  Operator learning, message passing, feature \& loss design, MESA/bbq, prior
  art \\
\ref{part:practice} & Working practice &
  Provenance labels, thresholds \& instruments, ADR switch conditions, guards
  that refuse, non-regenerable artifacts \\
\bottomrule
\end{tabularx}

\vspace{1em}
{\small The source document's own section list --- its original Part~0--VII
architecture in full --- is reproduced verbatim in
Appendix~\ref{app:sectionmap}.}

\cleardoublepage
\tableofcontents
\cleardoublepage
